\section{What It Serves, and What It Keeps}
\label{sec:ch08-what-it-serves-and-what-it-keeps}

\index{router!serves fewer types than the subgraph|(}%
Introspect the router, then introspect the subgraph behind it, and count the
types each one names. Discount the ones whose names begin with a double
underscore, which are introspection's own and belong to no schema. The router
serves 21. The subgraph serves 24.

\index{router!hides exactly the federation route}%
The three it does not pass on are these:

\begin{minted}{text}
_Any
_Entity
_Service
\end{minted}

\index{Query@\code{Query}!loses two fields at the router}%
And the same subtraction happens to the root. Ask each of them for the fields
on \code{Query} and the subgraph offers five against the router's three, or six
against four if you ask for deprecated fields as well. Either way the
difference is the same two: \code{\_entities} and \code{\_service}.
Chapter~\ref{ch:first-subgraph} built both, and
chapter~\ref{ch:composition} predicted from reading \code{graphqlSchema} inside
the config that they would be gone from what a client reads. They are. A client
talking to 3002 cannot see that federation is involved at all.

\subsection{Gone Is Not the Same as Closed}

\index{entities@\code{\_entities}!still open on the subgraph's own port}%
The subgraph did not lose those fields. It still has them, on 5001, and asking
it for its own schema still works:

\begin{minted}{graphql}
{ _service { sdl } }
\end{minted}

\begin{minted}[fontsize=\footnotesize,breakanywhere]{text}
{"data":{"_service":{"sdl":"schema\r\n  @link(\r\n    url: \"https://specs.apollo.dev/federation/v2.6\"\r\n    import: [\"@key\", \"@shareable\", \"@tag\", \"FieldSet\"]\r\n  ) {\r\n  query: Query\r\n ...
\end{minted}

\index{entities@\code{\_entities}!the router is not a control}%
That is the first 200 characters of one string, and the string runs to the
whole subgraph document: the \code{@link} header above, then every type, every
\code{@key} and the entity route itself, escaped newlines and all. It is
served to anyone who can reach the port.

\index{entities@\code{\_entities}!hidden is not closed}%
The router hiding a field is not the router closing it. It is one process
declining to advertise something that another process is still happy to answer,
and the only thing between the two is that nobody outside has a route to 5001
yet. On my laptop both ports are on the same machine and both are open.

I am stating that here and doing nothing about it, because
chapter~\ref{ch:entities-authorization} is entirely about what the entity route
will hand over when asked politely. What this chapter contributes is the shape
of the mistake: reading the router's schema and concluding that it describes
your attack surface. It describes one of your two attack surfaces.

\subsection{Two Things That Came Through Unchanged}

\index{FieldSet@\code{FieldSet}!is served to clients}%
Chapter~\ref{ch:composition} noticed a stray line in the composed client
schema, \code{scalar FieldSet}, left standing after the directive that needed
it was removed, and called it a generated document's untidiness. It is more than
a line in a file. Ask the router about it:

\begin{minted}{graphql}
{ __type(name: "FieldSet") { name kind } }
\end{minted}

\begin{minted}{json}
{"data":{"__type":{"name":"FieldSet","kind":"SCALAR"}}}
\end{minted}

\index{FieldSet@\code{FieldSet}!a scalar with no field of that type}%
So it is in the 21. Your clients are served a scalar that no field anywhere in
the schema is typed with, that they cannot send and cannot receive, and whose
only purpose belonged to a directive they will never see. Harmless. Also a
useful thing to have met once, for the day you diff two supergraph schemas and
find a type you cannot account for.

\index{deprecation!survives composition and the router}%
The other thing that survived is more useful.
Chapter~\ref{ch:schema-design} deprecated \code{sessionById} and wrote a reason
on it. Ask the router for its deprecated fields and the reason comes back word
for word:

\begin{minted}[breaklines]{text}
Ask for `node(id:)` instead. This field takes the database key, which is
only unique among sessions.
\end{minted}

\index{deprecation!is worth writing well because it travels}%
That sentence was written into a C\# attribute in chapter~\ref{ch:schema-design},
exported into a schema document, read by a composer, compiled into a router
execution config, and served to a client by a Go binary, without anybody
retyping it. Which is a reason to spend a minute on the wording: a deprecation
reason is one of the few pieces of English in your codebase that reaches your
consumers directly, and this is the path it takes to get there.
\index{router!serves fewer types than the subgraph|)}%
